% Architecture : couches, dépôt, outils, garde-fous (rapport de cadrage, F ; ADR-0022, 0025, 0028, 0030, 0040).

\chapter{Architecture}
\label{chap:architecture}

\section{Couches de données}

\begin{center}
\begin{tabular}{lll}
\toprule
Couche & Support & Produite par \\
\midrule
brut & fichiers \code{.json.gz} et CSV, journal & collecteurs \\
référentiel & Postgres, schéma \code{staging} & \code{load} (reconstruit) \\
variables & instantané Parquet versionné, \code{features.dataset\_version} & \code{features build} \\
modèles & fichiers \code{joblib}, cartes d'identité JSON & \code{modeling evaluate}, \code{train} \\
prédictions & Postgres, \code{ops.prediction} & inférence (API, commandes) \\
\bottomrule
\end{tabular}
\end{center}

Chaque couche se reconstruit depuis la précédente par une commande documentée. Le jeu de données est
un instantané : son identifiant (par exemple \code{ds-2026-10-22-bf786ef1}) porte la date et l'empreinte
de son manifeste (chargement lu, commit, paramètres, empreinte de chaque fichier), et une évaluation
cite toujours le jeu qu'elle a lu (ADR-0030).

\section{Code}

Le code Python est un paquet (\code{src/foot\_predictor/}) organisé par étape : \code{collect},
\code{ingestion}, \code{features}, \code{modeling}, \code{inference}, \code{api}, \code{ui}. Les
dépendances sont gérées par \code{uv} ; les migrations de la base par Alembic (8 révisions) ; les deux
bases (travail et test) tournent dans Docker. Les variables sont des \emph{fonctions pures} : elles
reçoivent un historique et rendent des valeurs, sans accès à la base, ce qui les rend testables sur
des données synthétiques. Chaque variable a sa fiche dans un registre (\code{features/registry.yaml}) :
définition, horizon, source.

\section{Garde-fous}

\begin{itemize}
  \item \textbf{Porte unique des données} : toute lecture des matchs passe par
        \code{features/sources.py} ; des tests d'architecture refusent tout autre accès (ADR-0028).
  \item \textbf{Règle temporelle} : une variable d'un match n'utilise que des données strictement
        antérieures au jour du match ; des tests d'\emph{invariance} modifient les matchs postérieurs et
        vérifient que rien ne bouge.
  \item \textbf{Aucun remplissage silencieux} : une valeur incalculable reste vide, et la ligne est
        exclue et comptée, jamais remplie par une moyenne.
  \item \textbf{Protocole figé} : le protocole d'évaluation est écrit et marqué d'un tag Git
        (\code{protocole-v1}) avant toute évaluation de modèle (ADR-0037).
  \item \textbf{Gel avant un test scellé} : la liste du test est figée par un tag (\code{pre-scelle-h1},
        \code{pre-scelle-h2}) ; jusqu'au test, le code d'évaluation et des variables ne change plus, et la
        commande refuse elle-même un code modifié depuis le tag ou un second lancement.
  \item \textbf{Garde du modèle testé} : une commande vérifie que le rejeu reproduit, à
        \(10^{-12}\) près, les prédictions des deux tests scellés (ADR-0052).
  \item \textbf{Secrets et données hors Git} : clés dans des fichiers \code{.env} ignorés, hooks de
        détection de secrets ; aucune donnée ni aucun modèle n'est versionné.
\end{itemize}

\section{Flux de travail}

Git suit un modèle simple (ADR-0022) : une seule branche longue, \code{main}, des branches courtes, une
demande de fusion (PR) par tâche, fusionnée après intégration continue verte (GitHub Actions : ruff et
957 tests), et des tags aux jalons (\code{v0.3.0} à \code{v1.0.0}, puis la version intermédiaire). La documentation suit la même
discipline : un état du projet mis à jour à chaque session, un mode d'emploi par étape, un retour écrit
à la fin de chaque lot de travail.

\section{Application}

L'inférence calcule les variables d'un match par la même fonction que l'entraînement. Une API FastAPI
l'expose en local (\code{127.0.0.1} seulement), et une interface Streamlit l'appelle en HTTP, sans jamais
importer la logique métier ni lire la base (ADR-0040). Le chapitre~\ref{chap:application} la décrit.
